% ============================================
% 3 模型假设与符号说明
% ============================================
\section{模型假设与符号说明}

\subsection{模型假设}

\begin{enumerate}[label=\arabic*.]
    \item \textbf{成分数据有效性假设}：各氧化物比例累加和在 $85\%\sim105\%$ 之间的采样数据视为有效，超出范围的样品不参与建模。
    \item \textbf{风化单向性假设}：风化过程是单向的（元素流失或富集有确定方向），风化后成分可以通过校正向量逆推风化前成分。
    \item \textbf{同类风化一致性假设}：同类型玻璃的风化效应具有统计一致性，可以用组间均值差或配对样本差作为校正向量。
    \item \textbf{未风化点代表性假设}：风化文物表面的"未风化点"可近似代表该文物风化前的原始成分。
    \item \textbf{成分闭合假设}：化学成分数据为成分性数据（总和为常数），在进行相关性分析时需通过 CLR 变换消除闭合效应。
    \item \textbf{分类独立性假设}：玻璃类型由其原始配方决定，不受风化影响——风化只改变成分比例，不改变类型归属。
\end{enumerate}

\subsection{符号说明}

\begin{table}[H]
    \centering
    \caption{符号说明表}
    \label{tab:symbols}
    \begin{tabular}{p{2cm} p{9cm} p{2cm}}
        \toprule
        \textbf{符号} & \textbf{含义} & \textbf{单位} \\
        \midrule
        $i$ & 文物/采样点索引 & —— \\
        $j$ & 氧化物索引（1--14） & —— \\
        $w_{ij}$ & 第 $i$ 个采样点第 $j$ 种氧化物的重量百分比 & \% \\
        $\Delta_j$ & 第 $j$ 种氧化物的风化效应（风化后 $-$ 风化前） & \% \\
        $C_{ij}^{\text{pre}}$ & 预测的风化前成分含量 & \% \\
        $C_{ij}^{\text{post}}$ & 实测的风化后成分含量 & \% \\
        $\text{CLR}(x)$ & centered log-ratio 变换 & —— \\
        $r_p$ & Pearson 相关系数 & —— \\
        $r_s$ & Spearman 秩相关系数 & —— \\
        $k$ & 聚类数（亚类个数） & —— \\
        $S(k)$ & 轮廓系数 & —— \\
        \bottomrule
    \end{tabular}
\end{table}
